\section{Analyse pratique de la heap}
\label{sec:Analyse_du_tas}

Vérifions ce que nous venons de dire précédemment.

Vous pouvez maintenant lancer ce programme avec pwndbg.
L'idée de ce code est d'observer les différentes "bin" et les métadonnées des chunks.

\begin{minted}[frame=single, framesep=2mm, fontsize=\footnotesize, style=monokai, bgcolor=pwndbgbg]{c}
#include <stdio.h>
#include <stdlib.h>
#include <string.h>

#define TCACHE_MAX 7

void fill_tcache(size_t size){
    void *ptrs[7];
    for (unsigned i = 0; i < TCACHE_MAX; i++)
        ptrs[i] = malloc(size);
    for (unsigned i = 0; i < TCACHE_MAX; i++)
        free(ptrs[i]);
}

int main(){

    // Tcache
    void *t1 = malloc(0x40), *t2 = malloc(0x100), *t3 = malloc(0x400);
    strcpy(t1, "T1T1T1T1"); // T = 0x54 et 1 = 0x31
    free(t1), free(t2),  free(t3);

    // Fastbin
    void *f_min = malloc(0x10),*g1 = malloc(0x30);
    void *f_max = malloc(0x70),*g2 = malloc(0x30);

    strcpy(f_max, "FFFFFFFFFFFFFFFFFFFFFFFF"); // F = 0x46

    fill_tcache(0x10); fill_tcache(0x70);

    free(f_min); free(f_max);

    // Smallbin
    void *s_min = malloc(0x80), *g3 = malloc(0x30);
    void *s_max = malloc(0x3e0), *g4 = malloc(0x30);

    strcpy(s_max, "SSSSSSSSSSSSSSSSSSSSSSSSSSSSSSSSSSSSSSSSSSSSSSSSSS"); // S = 0x53

    fill_tcache(0x80); fill_tcache(0x3e0);

    free(s_min); free(s_max);

    // Largebin
    void *l_min = malloc(0x3f0);

    fill_tcache(0x3f0);

    strcpy(l_min, "LLLLLLLLLLLLLLLLLLLLLLLLLLLLLLLLLLLLLLLLLLLLLLLL"); // L = 0x4c

    free(l_min);
    malloc(0x410);

    return 0;
}


\end{minted}


\section*{Code couleur utilisé pour l'analyse des chunks \texttt{glibc}}

Afin de faciliter la lecture et l'analyse des structures internes du tas
(\textit{heap}) dans \texttt{glibc}, chaque champ d'un chunk est représenté
avec une couleur spécifique.
Ce code couleur permet d'identifier rapidement le rôle de chaque octet
dans les dumps mémoire.

\begin{itemize}
    \item \textcolor{pwndbgbrown}{\textbf{Brun (\texttt{prev\_size})}} :
          Taille du chunk précédent. Ce champ n'est utilisé que lorsque le bit
          \texttt{PREV\_INUSE} est à zéro.

    \item \textcolor{pwndbgred}{\textbf{Rouge (\texttt{size})}} :
          Taille du chunk courant, incluant les bits d'état
          (\texttt{PREV\_INUSE}, \texttt{IS\_MMAPPED}, \texttt{NON\_MAIN\_ARENA}).

    \item \textcolor{pwndbgblue}{\textbf{Bleu (\texttt{fd})}} :
          Pointeur vers le prochain chunk dans une liste chaînée
          (fastbins, smallbins, unsorted bin, etc.).

    \item \textcolor{pwndbgpurple}{\textbf{Violet (\texttt{fd\_nextsize})}} :
          Pointeur utilisé dans les \textit{large bins} pour maintenir un tri
          secondaire par taille.

    \item \textcolor{pwndbgyellow}{\textbf{Jaune (\texttt{bk})}} :
          Pointeur vers le chunk précédent dans une liste chaînée.

    \item \textcolor{pwndbgorange}{\textbf{Orange (\texttt{bk\_nextsize})}} :
          Pointeur complémentaire à \texttt{fd\_nextsize} dans les \textit{large bins},
          permettant la navigation inverse dans la liste triée par taille.

    \item \textcolor{pwndbgdarkgreen}{\textbf{Vert foncé (\texttt{payload})}} :
          Données utilisateur stockées dans le chunk.
          C'est la zone retournée par \texttt{malloc()} et modifiable par le programme.


\end{itemize}

\subsection{Tcache}
\begin{minted}[frame=single, framesep=2mm, fontsize=\footnotesize, style=monokai, bgcolor=pwndbgbg]{c}
void *t1 = malloc(0x40), *t2 = malloc(0x100), *t3 = malloc(0x400);
strcpy(t1, "T1T1T1T1"); // T = 0x54 et 1 = 0x31
free(t1), free(t2), free(t3);
\end{minted}

Au début du programme, on n'a rien, pas de heap utilisateur.
Elle sera initialisée dès le premier \texttt{malloc}, la seule structure en place au début du programme est
la main arena qui est dans la libc. Vous pouvez voir son adresse avec \texttt{p \&main\_arena}.

Une fois que vous avez passé les trois \texttt{strcpy}, vous pouvez les retrouver avec la commande \texttt{heap}.
Vous allez ainsi voir 5 chunks, c'est normal.
Le premier chunk est un résidu de l'appel initial de \texttt{brk}, il sert à aligner le Top Chunk et à avoir
un premier chunk valide qui sert à la gestion des bins et à la recherche des chunks.

Voici à quoi ressemble t1 une fois free :

\begin{pwndbg}[x/40bx t1 - 0x10]{text}
    0x555555559290: !\textcolor{pwndbgbrown}{0x00}! !\textcolor{pwndbgbrown}{0x00}! !\textcolor{pwndbgbrown}{0x00}! !\textcolor{pwndbgbrown}{0x00}! !\textcolor{pwndbgbrown}{0x00}! !\textcolor{pwndbgbrown}{0x00}! !\textcolor{pwndbgbrown}{0x00}! !\textcolor{pwndbgbrown}{0x00}!
    0x555555559298: !\textcolor{pwndbgred}{0x51}! !\textcolor{pwndbgred}{0x00}! !\textcolor{pwndbgred}{0x00}! !\textcolor{pwndbgred}{0x00}! !\textcolor{pwndbgred}{0x00}! !\textcolor{pwndbgred}{0x00}! !\textcolor{pwndbgred}{0x00}! !\textcolor{pwndbgred}{0x00}!
    0x5555555592a0: !\textcolor{pwndbgblue}{0x59}! !\textcolor{pwndbgblue}{0x55}! !\textcolor{pwndbgblue}{0x55}! !\textcolor{pwndbgblue}{0x55}! !\textcolor{pwndbgblue}{0x05}! !\textcolor{pwndbgblue}{0x00}! !\textcolor{pwndbgblue}{0x00}! !\textcolor{pwndbgblue}{0x00}!
    0x5555555592a8: !\textcolor{pwndbgdarkgreen}{0x46}! !\textcolor{pwndbgdarkgreen}{0x01}! !\textcolor{pwndbgdarkgreen}{0x3a}! !\textcolor{pwndbgdarkgreen}{0x70}! !\textcolor{pwndbgdarkgreen}{0x5f}! !\textcolor{pwndbgdarkgreen}{0x9c}! !\textcolor{pwndbgdarkgreen}{0xce}! !\textcolor{pwndbgdarkgreen}{0x17}!
    0x5555555592b0: !\textcolor{pwndbgdarkgreen}{0x54}! !\textcolor{pwndbgdarkgreen}{0x31}! !\textcolor{pwndbgdarkgreen}{0x54}! !\textcolor{pwndbgdarkgreen}{0x31}! !\textcolor{pwndbgdarkgreen}{0x54}! !\textcolor{pwndbgdarkgreen}{0x31}! !\textcolor{pwndbgdarkgreen}{0x54}! !\textcolor{pwndbgdarkgreen}{0x31}!
\end{pwndbg}

Comme le montrent les couleurs, la zone utilisateur a été immédiatement écrasée par la GLIBC pour y stocker deux métadonnées critiques propres au Tcache :

En \textcolor{pwndbgblue}{bleu}, on retrouve le pointeur \texttt{fd} (forward) à la valeur 0x555555559 qui
indique l'adresse du prochain chunk libre à partir du champ utilisateur.
L'adresse semble illisible car, depuis la GLIBC 2.32, elle est protégée
par le \textit{Safe Linking}. Le pointeur n'est plus stocké en clair.
Il est masqué dynamiquement via une opération XOR.
Comme t1 est le premier chunk, son fd pointe sur 0.

$$
    \begin{aligned}
        \text{fd}_{\text{encodé}} & =  (Adresse\_Chunk \gg 12) \oplus \text{fd}_{\text{réel}} \\
                                  & = (0x5555555592a0 \gg 12) \oplus 0                        \\
                                  & =  0x555555559
    \end{aligned}
$$


En \textcolor{pwndbgdarkgreen}{vert foncé}, on observe le pointeur \texttt{key}.
Contrairement à ce que son nom suggère, ce n'est pas une clé servant à déchiffrer
le pointeur bleu. Il s'agit d'une sécurité anti-double-free introduite
dans la GLIBC 2.29, qui pointe vers la \texttt{tcache\_perthread\_struct}.
Si l'allocateur détecte cette signature lors d'un \texttt{free()},
il parcourt préventivement la liste chaînée pour s'assurer que le bloc n'y figure pas
déjà.

\subsection{Fastbin}
\begin{minted}[frame=single, framesep=2mm, fontsize=\footnotesize, style=monokai, bgcolor=pwndbgbg]{c}
      void *f_min = malloc(0x10),*g1 = malloc(0x30);
      void *f_max = malloc(0x70),*g2 = malloc(0x30);
      strcpy(f_max, "FFFFFFFFFFFFFFFFFFFFFFFF"); // F = 0x46
      fill_tcache(0x10); fill_tcache(0x70);
      free(f_min); free(f_max);
\end{minted}

Voici à quoi ressemble f\_max une fois free :

\begin{pwndbg}[x/40bx f\_max - 0x10]{text}
    0x555555559860: !\textcolor{pwndbgbrown}{0x00}! !\textcolor{pwndbgbrown}{0x00}! !\textcolor{pwndbgbrown}{0x00}! !\textcolor{pwndbgbrown}{0x00}! !\textcolor{pwndbgbrown}{0x00}! !\textcolor{pwndbgbrown}{0x00}! !\textcolor{pwndbgbrown}{0x00}! !\textcolor{pwndbgbrown}{0x00}!
    0x555555559868: !\textcolor{pwndbgred}{0x81}! !\textcolor{pwndbgred}{0x00}! !\textcolor{pwndbgred}{0x00}! !\textcolor{pwndbgred}{0x00}! !\textcolor{pwndbgred}{0x00}! !\textcolor{pwndbgred}{0x00}! !\textcolor{pwndbgred}{0x00}! !\textcolor{pwndbgred}{0x00}!
    0x555555559870: !\textcolor{pwndbgblue}{0x59}! !\textcolor{pwndbgblue}{0x55}! !\textcolor{pwndbgblue}{0x55}! !\textcolor{pwndbgblue}{0x55}! !\textcolor{pwndbgblue}{0x05}! !\textcolor{pwndbgblue}{0x00}! !\textcolor{pwndbgblue}{0x00}! !\textcolor{pwndbgblue}{0x00}!
    0x555555559878: !\textcolor{pwndbgdarkgreen}{0x46}! !\textcolor{pwndbgdarkgreen}{0x46}! !\textcolor{pwndbgdarkgreen}{0x46}! !\textcolor{pwndbgdarkgreen}{0x46}! !\textcolor{pwndbgdarkgreen}{0x46}! !\textcolor{pwndbgdarkgreen}{0x46}! !\textcolor{pwndbgdarkgreen}{0x46}! !\textcolor{pwndbgdarkgreen}{0x46}!
    0x555555559880: !\textcolor{pwndbgdarkgreen}{0x46}! !\textcolor{pwndbgdarkgreen}{0x46}! !\textcolor{pwndbgdarkgreen}{0x46}! !\textcolor{pwndbgdarkgreen}{0x46}! !\textcolor{pwndbgdarkgreen}{0x46}! !\textcolor{pwndbgdarkgreen}{0x46}! !\textcolor{pwndbgdarkgreen}{0x46}! !\textcolor{pwndbgdarkgreen}{0x46}!
\end{pwndbg}

Dans les fastbins, le \textit{Safe Linking} est actif.

Si l'on souhaite retrouver le \texttt{fd} réel, on peut manipuler la formule précédente
$$
    \begin{aligned}
        \text{fd}_{\text{réel}} & = \text{fd}_{\text{encodé}} \oplus (\text{adresse}_{\text{chunk}} \gg 12) \\
                                & = (0x555555559) \oplus (0x5555555592a0 \gg 12)                            \\
                                & = 0
    \end{aligned}
$$

Ici on trouve 0 car le chunk est seul dans le fastbin.


\subsection{Smallbin}
\begin{minted}[frame=single, framesep=2mm, fontsize=\footnotesize, style=monokai, bgcolor=pwndbgbg]{c}
      oid *s_min = malloc(0x80), *g3 = malloc(0x30);
      void *s_max = malloc(0x3e0), *g4 = malloc(0x30);
      strcpy(s_max, "SSSSSSSSSSSSSSSSSSSSSSSSSSSSSSSSSSSSSSSSSSSSSSSSSS"); // S = 0x53
      fill_tcache(0x80); fill_tcache(0x3e0);
      free(s_min); free(s_max);
\end{minted}


Voici à quoi ressemble s\_max une fois free :

\begin{pwndbg}[x/40bx s\_max - 0x10]{text}
    0x555555559e50: !\textcolor{pwndbgbrown}{0x00}! !\textcolor{pwndbgbrown}{0x00}! !\textcolor{pwndbgbrown}{0x00}! !\textcolor{pwndbgbrown}{0x00}! !\textcolor{pwndbgbrown}{0x00}! !\textcolor{pwndbgbrown}{0x00}! !\textcolor{pwndbgbrown}{0x00}! !\textcolor{pwndbgbrown}{0x00}!
    0x555555559e58: !\textcolor{pwndbgred}{0xf1}! !\textcolor{pwndbgred}{0x03}! !\textcolor{pwndbgred}{0x00}! !\textcolor{pwndbgred}{0x00}! !\textcolor{pwndbgred}{0x00}! !\textcolor{pwndbgred}{0x00}! !\textcolor{pwndbgred}{0x00}! !\textcolor{pwndbgred}{0x00}!
    0x555555559e60: !\textcolor{pwndbgblue}{0x80}! !\textcolor{pwndbgblue}{0x9d}! !\textcolor{pwndbgblue}{0x55}! !\textcolor{pwndbgblue}{0x55}! !\textcolor{pwndbgblue}{0x55}! !\textcolor{pwndbgblue}{0x55}! !\textcolor{pwndbgblue}{0x00}! !\textcolor{pwndbgblue}{0x00}!
    0x555555559e68: !\textcolor{pwndbgyellow}{0x20}! !\textcolor{pwndbgyellow}{0x3b}! !\textcolor{pwndbgyellow}{0xe0}! !\textcolor{pwndbgyellow}{0xf7}! !\textcolor{pwndbgyellow}{0xff}! !\textcolor{pwndbgyellow}{0x7f}! !\textcolor{pwndbgyellow}{0x00}! !\textcolor{pwndbgyellow}{0x00}!
    0x555555559e70: !\textcolor{pwndbgdarkgreen}{0x53}! !\textcolor{pwndbgdarkgreen}{0x53}! !\textcolor{pwndbgdarkgreen}{0x53}! !\textcolor{pwndbgdarkgreen}{0x53}! !\textcolor{pwndbgdarkgreen}{0x53}! !\textcolor{pwndbgdarkgreen}{0x53}! !\textcolor{pwndbgdarkgreen}{0x53}! !\textcolor{pwndbgdarkgreen}{0x53}!
\end{pwndbg}


Si l’on regarde l’état des bins, on s’attendrait à avoir un chunk dans l’unsorted bin
(\texttt{s\_max}) et un dans le smallbin (\texttt{s\_min}), puisque \texttt{s\_max}
est plus grand que \texttt{s\_min}. Cependant, ce n’est pas le cas, car un appel à
\texttt{free} ne déclenche jamais le tri des chunks~: il effectue uniquement une éventuelle
consolidation avec les chunks voisins.


\begin{pwndbg}[bins]{text}
    tcachebins
    0x20 [  7]: 0x5555555599f0 -> .... <-0
    0x50 [  1]: t1 <-0
    0x80 [  7]: 0x555555559d10 -> ....  <-0
    0x90 [  7]: 0x55555555a5f0 -> ....  <-0
    0x110 [  1]: t2 <-0
    0x3f0 [  7]: 0x55555555be20 -> .... <-0
    0x410 [  1]: t3 <-0
    fastbins
    0x20: f_min <-0
    0x80: f_max <-0
    unsortedbin
    all: s_max -> s_min -> 0x7ffff7e03b20 (main_arena+96) <- s_max
    smallbins
    empty
    largebins
    empty

\end{pwndbg}



\subsection{Largebin}
\begin{minted}[frame=single, framesep=2mm, fontsize=\footnotesize, style=monokai, bgcolor=pwndbgbg]{c}
    void *l_min = malloc(0x3f0);
    fill_tcache(0x3f0);
    strcpy(l_min, "LLLLLLLLLLLLLLLLLLLLLLLLLLLLLLLLLLLLLLLLLLLLLLLL"); // L = 0x4c
    free(l_min);
    malloc(0x410);
\end{minted}

Lors de l'allocation d'un grand chunk, l'allocateur parcourt les différents bins dans
l'espoir de trouver des chunks qu'il pourrait fusionner.
Tous les chunks parcourus sont ensuite directement triés et se retrouvent soit dans un smallbin, soit dans
un largebin.

Voici l'état de \texttt{l\_min} après libération :

\begin{pwndbg}[x/56bx l\_min - 0x10]{text}
    0x55555555c200: !\textcolor{pwndbgbrown}{0x00}! !\textcolor{pwndbgbrown}{0x00}! !\textcolor{pwndbgbrown}{0x00}! !\textcolor{pwndbgbrown}{0x00}! !\textcolor{pwndbgbrown}{0x00}! !\textcolor{pwndbgbrown}{0x00}! !\textcolor{pwndbgbrown}{0x00}! !\textcolor{pwndbgbrown}{0x00}!
    0x55555555c208: !\textcolor{pwndbgred}{0x01}! !\textcolor{pwndbgred}{0x04}! !\textcolor{pwndbgred}{0x00}! !\textcolor{pwndbgred}{0x00}! !\textcolor{pwndbgred}{0x00}! !\textcolor{pwndbgred}{0x00}! !\textcolor{pwndbgred}{0x00}! !\textcolor{pwndbgred}{0x00}!
    0x55555555c210: !\textcolor{pwndbgblue}{0x10}! !\textcolor{pwndbgblue}{0x3f}! !\textcolor{pwndbgblue}{0xe0}! !\textcolor{pwndbgblue}{0xf7}! !\textcolor{pwndbgblue}{0xff}! !\textcolor{pwndbgblue}{0x7f}! !\textcolor{pwndbgblue}{0x00}! !\textcolor{pwndbgblue}{0x00}!
    0x55555555c218: !\textcolor{pwndbgyellow}{0x10}! !\textcolor{pwndbgyellow}{0x3f}! !\textcolor{pwndbgyellow}{0xe0}! !\textcolor{pwndbgyellow}{0xf7}! !\textcolor{pwndbgyellow}{0xff}! !\textcolor{pwndbgyellow}{0x7f}! !\textcolor{pwndbgyellow}{0x00}! !\textcolor{pwndbgyellow}{0x00}!
    0x55555555c220: !\textcolor{pwndbgpurple}{0x00}! !\textcolor{pwndbgpurple}{0xc2}! !\textcolor{pwndbgpurple}{0x55}! !\textcolor{pwndbgpurple}{0x55}! !\textcolor{pwndbgpurple}{0x55}! !\textcolor{pwndbgpurple}{0x55}! !\textcolor{pwndbgpurple}{0x00}! !\textcolor{pwndbgpurple}{0x00}!
    0x55555555c228: !\textcolor{pwndbgorange}{0x00}! !\textcolor{pwndbgorange}{0xc2}! !\textcolor{pwndbgorange}{0x55}! !\textcolor{pwndbgorange}{0x55}! !\textcolor{pwndbgorange}{0x55}! !\textcolor{pwndbgorange}{0x55}! !\textcolor{pwndbgorange}{0x00}! !\textcolor{pwndbgorange}{0x00}!
    0x55555555c230: !\textcolor{pwndbgdarkgreen}{0x4c}! !\textcolor{pwndbgdarkgreen}{0x4c}! !\textcolor{pwndbgdarkgreen}{0x4c}! !\textcolor{pwndbgdarkgreen}{0x4c}! !\textcolor{pwndbgdarkgreen}{0x4c}! !\textcolor{pwndbgdarkgreen}{0x4c}! !\textcolor{pwndbgdarkgreen}{0x4c}! !\textcolor{pwndbgdarkgreen}{0x4c}!
\end{pwndbg}


on observe que les pointeurs
\textcolor{pwndbgblue}{fd} et \textcolor{pwndbgyellow}{bk} pointent tous deux vers la tête du
largebin, ce qui est normal pour un bin circulaire contenant un seul élément.
Les champs \textcolor{pwndbgpurple}{fd\_nextsize} et \textcolor{pwndbgorange}{bk\_nextsize}
pointent vers lui-même, car il est le seul chunk de cette classe de taille.

On remarque que les deux chunks qui se trouvaient dans les fastbins se retrouvent
désormais dans les smallbins. Cela est dû à la consolidation déclenchée par
\texttt{malloc}, qui a vidé les fastbins et a directement rangé ces chunks dans
les smallbins correspondants.


\begin{pwndbg}[bins]{text}
    tcachebins
    0x20 [  7]: 0x5555555599f0 -> .... <-0
    0x50 [  1]: t1 <-0
    0x80 [  7]: 0x555555559d10 -> ....  <-0
    0x90 [  7]: 0x55555555a5f0 -> ....  <-0
    0x110 [  1]: t2 <-0
    0x3f0 [  7]: 0x55555555be20 -> .... <-0
    0x410 [  1]: t3 <-0
    fastbins
    empty
    unsortedbin
    empty
    smallbins
    0x20: f_min -> 0x7ffff7e03b30 (main_arena+112) <- f_min
    0x80: f_max -> 0x7ffff7e03b90 (main_arena+208) <- f_max
    0x90: s_min -> 0x7ffff7e03ba0 (main_arena+224) <- s_min
    0x3f0: s_max -> 0x7ffff7e03f00 (main_arena+1088) <- s_max
    largebins
    0x400-0x430: l_min -> 0x7ffff7e03f10 (main_arena+1104) <- l_min
\end{pwndbg}

